\documentclass[11pt,a4paper,oneside]{article}

% -------------------------------------------------
% Basic packages
% -------------------------------------------------
\usepackage[T1]{fontenc}
\usepackage[utf8]{inputenc}
\usepackage[spanish,es-noquoting,es-nodecimaldot]{babel}
\usepackage{lmodern}
\usepackage{microtype}
\usepackage[a4paper,margin=1.18in,headheight=15pt]{geometry}
\usepackage{amsmath,amssymb,amsfonts,amsthm,mathtools}
\usepackage{bm}
\usepackage{enumitem}
\usepackage{booktabs}
\usepackage{array}
\usepackage{xcolor}
\usepackage{hyperref}
\usepackage{setspace}
\usepackage{fancyhdr}
\usepackage{titlesec}
\usepackage[most]{tcolorbox}
\usepackage[capitalize,noabbrev]{cleveref}

\hypersetup{
  colorlinks=true,
  linkcolor=blue!50!black,
  citecolor=red!55!black,
  urlcolor=teal!60!black,
  pdftitle={Fundamentos de Econometría -- Tema 1: Principios de Econometría},
  pdfauthor={César Galindo}
}

% -------------------------------------------------
% Theorem-like environments (usados como "Definición"/"Ejemplo"/"Ejercicio")
% -------------------------------------------------
\theoremstyle{definition}
\newtheorem{definition}{Definición}[section]
\newtheorem{example}[definition]{Ejemplo}
\newtheorem{exercise}[definition]{Ejercicio}

% -------------------------------------------------
% Fancy page and heading style
% -------------------------------------------------
\definecolor{ihesblue}{RGB}{25,62,114}
\definecolor{ihesgold}{RGB}{160,120,40}
\pagestyle{fancy}
\fancyhf{}
\fancyhead[L]{\textsc{Fundamentos de Econometría -- ECO515}}
\fancyhead[R]{\textsc{César Galindo}}
\fancyfoot[C]{\thepage}
\renewcommand{\headrulewidth}{0.35pt}
\renewcommand{\footrulewidth}{0pt}
\titleformat{\section}{\Large\bfseries\color{ihesblue}}{\thesection.}{0.65em}{}
\titleformat{\subsection}{\large\bfseries\color{ihesblue}}{\thesubsection.}{0.65em}{}
\titlespacing*{\section}{0pt}{1.9em}{1.05em}
\titlespacing*{\subsection}{0pt}{1.45em}{0.95em}
\tcbset{colback=white,colframe=ihesblue!65!black,boxrule=0.5pt,arc=1.5pt}
\setlength{\parskip}{0.85em}
\setstretch{1.13}
\setlength{\parindent}{0pt}

% -------------------------------------------------
% Pedagogical boxes
% -------------------------------------------------
\newtcolorbox{ideaBox}[1][]{enhanced,breakable,colback=blue!2,colframe=ihesblue,
  borderline west={1.5pt}{0pt}{ihesblue},title=Idea central,fonttitle=\bfseries,#1}
\newtcolorbox{summaryBox}[1][]{enhanced,breakable,colback=gray!4,colframe=black!55,
  borderline west={1.5pt}{0pt}{black!70},title=Qué conviene recordar,fonttitle=\bfseries,#1}
\newtcolorbox{warningBox}[1][]{enhanced,breakable,colback=red!2,colframe=red!55!black,
  borderline west={1.5pt}{0pt}{red!75!black},title=Error frecuente,fonttitle=\bfseries,#1}
\newtcolorbox{empresaBox}[1][]{enhanced,breakable,colback=orange!3!white,colframe=orange!60!black,
  borderline west={1.5pt}{0pt}{orange!70!black},title=Conexión empresarial,fonttitle=\bfseries\color{orange!60!black},#1}
\newtcolorbox{cfaBox}[1][]{enhanced,breakable,
  colback=violet!3!white,colframe=violet!50!black,
  borderline west={2pt}{0pt}{violet!65!black},
  title=Ejercicio estilo CFA,fonttitle=\bfseries\color{violet!80!black},
  attach boxed title to top left={yshift=-2mm,xshift=4mm},
  boxed title style={colback=violet!12,colframe=violet!50!black,boxrule=0.5pt},#1}
\newtcolorbox{practiceBox}[1][]{enhanced,breakable,
  colback=green!2!white,colframe=green!45!black,
  borderline west={2pt}{0pt}{green!55!black},
  title=Para practicar,fonttitle=\bfseries\color{green!50!black},
  attach boxed title to top left={yshift=-2mm,xshift=4mm},
  boxed title style={colback=green!10,colframe=green!45!black,boxrule=0.5pt},#1}

\title{\vspace{-1.5cm}\textbf{\Large Fundamentos de Econometría}\\[2pt]\large Tema 1: Principios de Econometría}
\author{César Galindo\\ \small Universidad Panamericana \textperiodcentered\ Licenciatura en Finanzas \textperiodcentered\ ECO515}
\date{\small Otoño 2026}

\begin{document}
\maketitle

\begin{ideaBox}
Este documento cubre el Tema 1 del syllabus de \emph{Fundamentos de Econometría}
(ECO515, Semana 1): qué es la econometría, cómo se distingue la econometría
financiera de la económica, de qué piezas se compone un modelo econométrico,
cómo se clasifican los datos económicos y financieros, y de dónde viene la
disciplina. Incluye ejercicios de práctica y una sección de ejercicios estilo
CFA (Quantitative Methods, Nivel I) sobre estos mismos conceptos.
\end{ideaBox}

\tableofcontents

\section{¿Qué es la econometría?}

La \textbf{econometría} es la disciplina que usa la teoría económica, las
matemáticas y la inferencia estadística para cuantificar relaciones
económicas, poner a prueba hipótesis y pronosticar variables económicas y
financieras a partir de datos reales.

\begin{definition}[Econometría]
La econometría es el análisis cuantitativo de fenómenos económicos reales,
basado en el desarrollo simultáneo de la teoría y la observación, relacionadas
mediante métodos apropiados de inferencia (definición atribuida a Ragnar
Frisch, cofundador de la \emph{Econometric Society} en 1930).
\end{definition}

La palabra misma lo dice: \emph{econo-metría} = medición de la economía. Pero
medir no es solo calcular un promedio: es traducir una afirmación cualitativa
de la teoría (``si sube la tasa de interés, baja la inversión'') en una
relación cuantificable, con un número y un margen de error asociados
(``un aumento de 1 punto porcentual en la tasa reduce la inversión en
$\beta$ millones de pesos, con un intervalo de confianza de...'').

\begin{summaryBox}
La econometría combina tres ingredientes, y ninguno basta por sí solo:
\begin{itemize}[nosep]
\item \textbf{Teoría económica o financiera:} qué relación se espera y por qué (el signo esperado de $\beta$, la forma funcional).
\item \textbf{Matemáticas:} el modelo formal (ecuaciones, álgebra matricial, cálculo de optimización para estimar parámetros).
\item \textbf{Estadística (inferencia):} cómo estimar los parámetros a partir de una muestra y cuantificar la incertidumbre alrededor de esa estimación.
\end{itemize}
\end{summaryBox}

\subsection{Econometría descriptiva vs.\ econometría teórica}
Es útil distinguir, dentro de la econometría misma, entre:
\begin{itemize}[nosep]
\item \textbf{Econometría teórica:} desarrolla los métodos mismos (propiedades
de los estimadores, condiciones bajo las cuales MCO es insesgado o eficiente,
etc.). No es el foco de este curso, aunque se usarán sus resultados.
\item \textbf{Econometría aplicada:} usa esos métodos para estimar modelos
sobre datos reales y responder preguntas económicas o financieras concretas.
Éste es el foco del curso.
\end{itemize}

\section{Econometría para Economía vs.\ Econometría Financiera}

Aunque comparten las mismas herramientas estadísticas (regresión, pruebas de
hipótesis, series de tiempo), la \textbf{econometría económica} y la
\textbf{econometría financiera} difieren en el tipo de preguntas que hacen y,
por lo tanto, en qué supuestos son razonables.

\begin{center}
\renewcommand{\arraystretch}{1.35}
\begin{tabular}{p{4.1cm}p{5.2cm}p{5.2cm}}
\toprule
& \textbf{Econometría económica} & \textbf{Econometría financiera} \\
\midrule
\textbf{Preguntas típicas} & ¿Cómo afecta el gasto público al PIB? ¿Cuál es la elasticidad-precio de la demanda? & ¿El mercado es eficiente? ¿Qué explica el rendimiento de un activo? ¿Cuál es su volatilidad? \\
\textbf{Frecuencia de los datos} & Trimestral o anual, series relativamente cortas & Diaria, incluso intradía; series largas y muy ruidosas \\
\textbf{Supuestos típicos que fallan} & Endogeneidad, variables omitidas & Heteroscedasticidad (volatilidad cambiante), colas pesadas, no normalidad \\
\textbf{Modelos característicos} & Regresión múltiple, modelos de panel, modelos macro estructurales & CAPM y modelos multifactoriales, GARCH, modelos de series de tiempo financieras \\
\bottomrule
\end{tabular}
\end{center}

\begin{empresaBox}
El ejemplo canónico de econometría financiera es el \textbf{CAPM}:
$E(R_i) - R_f = \beta_i\,[E(R_m)-R_f]$. Estimar $\beta_i$ por MCO, usando
rendimientos históricos del activo $i$ y del mercado, es literalmente
correr una regresión lineal simple --- el mismo método que se introduce en
el Tema 2 de este curso, aplicado a un problema de finanzas en vez de uno
macroeconómico.
\end{empresaBox}

\section{Elementos de un modelo econométrico}

Todo modelo econométrico simple comparte la misma anatomía. Considérese la
forma general de un modelo de regresión lineal:
\[
Y_i = \beta_0 + \beta_1 X_i + u_i.
\]

\begin{definition}[Elementos de un modelo econométrico]
\begin{itemize}[nosep]
\item \textbf{Variable dependiente (o explicada), $Y$:} lo que se quiere explicar o pronosticar (p.\,ej., el rendimiento de una acción).
\item \textbf{Variable(s) independiente(s) (o explicativa(s)), $X$:} lo que se usa para explicar a $Y$ (p.\,ej., el rendimiento del mercado).
\item \textbf{Parámetros, $\beta_0,\beta_1$:} números fijos, desconocidos, que se estiman con los datos; cuantifican la relación entre $X$ y $Y$.
\item \textbf{Término de error (o perturbación), $u$:} recoge todo lo que afecta a $Y$ y no está en el modelo --- variables omitidas, error de medición, aleatoriedad genuina.
\item \textbf{Forma funcional:} cómo se relacionan $Y$ y $X$ (lineal, log-lineal, cuadrática, \dots); el Tema 2 del curso trata este punto con más detalle.
\end{itemize}
\end{definition}

\begin{example}
En el modelo $\text{Rendimiento}_i = \beta_0 + \beta_1\,\text{Rendimiento del mercado}_i + u_i$:
$Y$ es el rendimiento de la acción $i$; $X$ es el rendimiento del mercado;
$\beta_1$ (la beta del CAPM) mide la sensibilidad de la acción al mercado;
$\beta_0$ es el rendimiento esperado cuando el mercado no se mueve; y $u_i$
captura el riesgo idiosincrático (específico de la empresa) no explicado por
el mercado.
\end{example}

\begin{warningBox}
El término de error $u_i$ \textbf{no} es ``el error del analista''; es una
parte estructural del modelo. Ningún modelo econométrico pretende explicar
el 100\,\% de la variación de $Y$: si $u_i\equiv0$ para todos los datos, algo
está mal (sobreajuste, o $Y$ y $X$ son en realidad la misma variable
disfrazada).
\end{warningBox}

\section{Tipos y formas de agrupación de datos}

Antes de estimar cualquier modelo hay que saber qué tipo de datos se tiene,
porque eso determina qué métodos son válidos.

\begin{definition}[Tipos de datos económicos y financieros]
\begin{itemize}[nosep]
\item \textbf{Series de tiempo (\emph{time series}):} observaciones de \emph{una}
unidad (un país, una empresa, un activo) a lo largo del tiempo, en orden
cronológico. Ejemplo: el precio diario de cierre del IPC durante 2020--2026.
\item \textbf{Datos transversales (\emph{cross-section}):} observaciones de
\emph{muchas} unidades en \emph{un solo} punto en el tiempo. Ejemplo: el
rendimiento anual de las 35 empresas de la muestra del IPC en 2025.
\item \textbf{Datos de panel (\emph{panel data}):} combinan ambas dimensiones ---
muchas unidades, observadas a lo largo de varios periodos. Ejemplo: el
rendimiento trimestral de 35 empresas del IPC durante 20 trimestres.
\item \textbf{Datos continuos vs.\ discretos:} una variable es \emph{continua}
si puede tomar cualquier valor real en un intervalo (un precio, una tasa de
rendimiento); es \emph{discreta} si solo toma valores en un conjunto
numerable, frecuentemente enteros (el número de transacciones en un día, una
variable binaria como ``la empresa pagó dividendos: sí/no'').
\end{itemize}
\end{definition}

\begin{summaryBox}
Un panel ``ve'' tanto la variación \emph{entre} unidades (¿por qué la
empresa A rinde más que la B, en promedio?) como la variación \emph{dentro}
de cada unidad a través del tiempo (¿por qué la empresa A rindió más este
año que el pasado?). Una serie de tiempo pura solo ve la segunda; un corte
transversal puro solo ve la primera.
\end{summaryBox}

\begin{practiceBox}
Clasifica cada conjunto de datos como serie de tiempo, corte transversal, o
panel, y di si la variable de interés es continua o discreta:
\begin{enumerate}[label=(\alph*)]
\item El PIB trimestral de México, 2000--2026.
\item La calificación crediticia (AAA, AA, \dots, D) de 50 empresas emisoras de deuda, al cierre de 2025.
\item El rendimiento mensual de 20 fondos de inversión, durante 36 meses.
\item El número de operaciones (\emph{trades}) ejecutadas cada minuto en una acción durante un día de bolsa.
\end{enumerate}
\emph{Respuestas:} (a) serie de tiempo, continua; (b) corte transversal,
discreta (categórica ordinal); (c) panel, continua; (d) serie de tiempo (alta
frecuencia), discreta (conteo).
\end{practiceBox}

\section{Definición, antecedentes, objeto y método de la econometría}

\subsection{Antecedentes históricos}
El término ``econometría'' (del francés \emph{économétrie}) fue acuñado
formalmente con la fundación de la \textbf{Econometric Society} en 1930, por
iniciativa de Ragnar Frisch, Charles Roos e Irving Fisher. Frisch y Jan
Tinbergen recibieron el primer Premio Nobel de Economía en 1969
``por haber desarrollado y aplicado modelos dinámicos al análisis de
procesos económicos''. Trygve Haavelmo, en su artículo de 1944 \emph{The
Probability Approach in Econometrics}, sentó las bases probabilísticas
modernas de la disciplina: tratar los datos económicos como realizaciones de
variables aleatorias, no como cantidades determinísticas --- exactamente el
punto de vista que permite hablar de errores estándar e intervalos de
confianza para $\beta_1$.

\subsection{Objeto de estudio}
El objeto de la econometría es doble:
\begin{enumerate}[nosep]
\item \textbf{Estimar} relaciones cuantitativas entre variables económicas o financieras a partir de datos observados (no experimentales, en la gran mayoría de los casos).
\item \textbf{Contrastar} (poner a prueba) predicciones de la teoría económica o financiera contra la evidencia empírica.
\end{enumerate}

\subsection{Método}
El método econométrico típico sigue estos pasos, que reaparecerán a lo largo
del curso:
\begin{enumerate}[nosep]
\item Plantear una hipótesis teórica (p.\,ej., el CAPM predice $\beta_1>0$).
\item Especificar un modelo matemático que la traduzca en una ecuación estimable.
\item Recolectar los datos relevantes y clasificarlos (Sección 4).
\item Estimar los parámetros (Tema 2: mínimos cuadrados ordinarios).
\item Poner a prueba la hipótesis mediante inferencia estadística (Tema 3).
\item Diagnosticar el modelo (Temas 5--6) y, si es necesario, corregirlo o re-especificarlo.
\item Usar el modelo para pronosticar o para política/decisión de inversión.
\end{enumerate}

\section{Ejercicios}

\begin{exercise}
Enuncia, con tus propias palabras, la diferencia entre econometría teórica y
econometría aplicada. Da un ejemplo de cada una relacionado con mercados
financieros.
\end{exercise}

\begin{exercise}
Para el modelo $\text{Precio de la vivienda}_i = \beta_0 + \beta_1\,\text{Metros cuadrados}_i + u_i$,
identifica la variable dependiente, la variable independiente, los
parámetros y el término de error. ¿Qué signo esperarías para $\beta_1$ y por
qué?
\end{exercise}

\begin{exercise}
Un analista tiene el rendimiento diario de las 35 acciones del IPC durante
los últimos 5 años. Explica por qué estos datos son un panel, y describe
brevemente una pregunta que solo se puede responder aprovechando la
dimensión de panel (y no con un simple corte transversal o una simple serie
de tiempo).
\end{exercise}

\begin{exercise}
¿Por qué, según Haavelmo (1944), es necesario un enfoque probabilístico para
poder hablar de ``errores estándar'' o ``intervalos de confianza'' en
econometría? ¿Qué pasaría con esos conceptos si tratáramos los datos
económicos como cantidades puramente determinísticas?
\end{exercise}

\section{Ejercicios estilo CFA}

Las siguientes preguntas siguen el formato de opción múltiple (A/B/C) del
programa CFA (Nivel I, \emph{Quantitative Methods}), aplicadas a los
conceptos de este tema.

\begin{cfaBox}
\textbf{1.} Un analista recopila el rendimiento trimestral de 40 fondos de
cobertura (\emph{hedge funds}) distintos, cada uno observado durante 12
trimestres consecutivos. Este conjunto de datos se describe \textbf{mejor}
como:
\begin{enumerate}[label=\Alph*.]
\item Una serie de tiempo.
\item Un corte transversal.
\item Un panel de datos.
\end{enumerate}
\end{cfaBox}
\emph{Solución: C.} Hay múltiples unidades (40 fondos) observadas a través
de múltiples periodos (12 trimestres) para cada una: es la combinación
característica de un panel. (A) sería correcto solo si hubiera un único
fondo; (B) sería correcto solo si cada fondo se observara en un único
trimestre.

\begin{cfaBox}
\textbf{2.} En el modelo de regresión $\text{Rendimiento}_{i} = \alpha + \beta\,(\text{Rendimiento del mercado}_i) + \varepsilon_i$,
el coeficiente $\beta$ representa:
\begin{enumerate}[label=\Alph*.]
\item El rendimiento esperado del activo cuando el mercado no se mueve.
\item La sensibilidad del rendimiento del activo ante cambios en el rendimiento del mercado.
\item La parte del rendimiento del activo que no explica el mercado.
\end{enumerate}
\end{cfaBox}
\emph{Solución: B.} $\beta$ es el parámetro de pendiente: mide cuánto cambia
$Y$ (el rendimiento del activo) por cada unidad de cambio en $X$ (el
rendimiento del mercado). (A) describe al intercepto $\alpha$; (C) describe
al término de error $\varepsilon_i$.

\begin{cfaBox}
\textbf{3.} ¿Cuál de las siguientes variables es \textbf{discreta}?
\begin{enumerate}[label=\Alph*.]
\item El precio de cierre diario de una acción.
\item El número de operaciones de compra ejecutadas en un fondo durante un mes.
\item La tasa interna de retorno (TIR) de un proyecto de inversión.
\end{enumerate}
\end{cfaBox}
\emph{Solución: B.} Un conteo de operaciones solo puede tomar valores
enteros no negativos ($0,1,2,\dots$): es discreta. Tanto un precio como una
TIR pueden, en principio, tomar cualquier valor real dentro de un rango:
ambas son continuas.

\begin{cfaBox}
\textbf{4.} Según el enfoque probabilístico de Haavelmo (1944), los datos
económicos observados se interpretan \textbf{principalmente} como:
\begin{enumerate}[label=\Alph*.]
\item Cantidades determinísticas, libres de todo error de medición.
\item Realizaciones particulares de variables aleatorias subyacentes.
\item Parámetros poblacionales conocidos con certeza.
\end{enumerate}
\end{cfaBox}
\emph{Solución: B.} El aporte central de Haavelmo fue tratar los datos
económicos observados como una muestra (una realización) de un proceso
aleatorio subyacente, lo cual es precisamente lo que permite aplicar
inferencia estadística --- errores estándar, pruebas de hipótesis,
intervalos de confianza --- a los parámetros estimados.

\begin{cfaBox}
\textbf{5.} Un analista de econometría financiera, a diferencia de uno de
econometría económica ``pura'', debe prestar \textbf{particular} atención a
cuál de los siguientes problemas, dado que trabaja frecuentemente con series
de rendimientos financieros de alta frecuencia:
\begin{enumerate}[label=\Alph*.]
\item Estacionalidad en el gasto público trimestral.
\item Heteroscedasticidad (volatilidad cambiante en el tiempo).
\item Elasticidades de largo plazo en modelos de demanda agregada.
\end{enumerate}
\end{cfaBox}
\emph{Solución: B.} Los rendimientos financieros exhiben típicamente
``agrupamiento de volatilidad'' (\emph{volatility clustering}): periodos de
alta y baja varianza que se alternan --- la razón por la que modelos como
GARCH (fuera del alcance de este tema, pero relevante para econometría
financiera) son tan usados en finanzas. (A) y (C) son preocupaciones más
características de la econometría macroeconómica.

\begin{summaryBox}
\textbf{En una frase:} la econometría traduce afirmaciones cualitativas de
la teoría económica o financiera en modelos cuantificables y contrastables
--- y el primer paso, antes de estimar nada, es reconocer qué tipo de datos
se tiene y qué piezas tendrá el modelo.
\end{summaryBox}

\end{document}
